\documentclass[11pt]{article}
\usepackage{fontspec,amsmath,amssymb,amsthm,hyperref}
\usepackage[margin=1in]{geometry}
\setlength{\emergencystretch}{3em}
\title{Construct the empirical parameter Brownian input}
\author{LangevinLean mathematical exposition}
\date{}
\begin{document}\maketitle
\noindent Source module: \texttt{ProximalLangevin/MeanFieldNoise.lean}. Compile this proof document with XeLaTeX or LuaLaTeX. The accompanying HTML places the complete Lean declarations beside every proof.


On \(Q^{\otimes N}\), \(N>0\), use copy zero of a supplied \(d_\theta\)-coordinate Brownian source for the common parameter input. The actual noise amplitude is \(\sqrt{2/N}\).


\section{Pull back Brownian coordinates and define their scaled L\ensuremath{{}^{2}} noise}


Coordinate evaluation on the finite product is measure preserving, so each pulled-back scalar process retains its standard Brownian law. Define \(A_t=N^{-1/2}\sqrt2(B_t^\theta-B_0^\theta)\) as an actual \(L^2\) element. The Brownian-noise construction proves continuity in \(L^2\) and value zero at time zero; scalar multiplication preserves both facts.



\noindent Lean declarations: \texttt{empiricalParameter\_preBrownian}, \texttt{empiricalParameterNoise}, \texttt{empiricalParameterNoise\_continuous}, \texttt{empiricalParameterNoise\_zero}.

\section{Derive the exact norm and its N-independent horizon coefficient}


The coordinate second moments sum to \(d_\theta t\). Consequently

\[\|A_t\|_2=\frac{\sqrt{2d_\theta}\sqrt t}{\sqrt N}\le\frac{\sqrt{2d_\theta}\sqrt T}{\sqrt N},\qquad0\le t\le T.\]

The theorem follows by the exact Brownian norm identity and monotonicity of the square root. The canonical specialization uses the already constructed Gaussian time-indicator Brownian source, so it supplies a concrete instance rather than postulating one.



\noindent Lean declarations: \texttt{empiricalParameterNoise\_norm}, \texttt{empiricalParameterNoise\_bound}, \texttt{canonical\_empiricalParameterNoise\_bound}.

\end{document}
